\section{Step 5: Blueprint Planning via a Three-Stage Stochastic Program}

OR scheduling problems decompose into three planning levels~\cite{samudra2016}:
the Case Mix Problem (CMP) at the strategic level, the Master Surgery
Scheduling Problem (MSSP) at the tactical level, and the Surgery Scheduling
Problem (SSP) at the operational level.
The Step~3 model addresses the SSP; this step adds a tactical MSSP-level
first stage, yielding a three-stage stochastic program whose new first stage
decides how many patients of each type to schedule per session before the
waiting list is known, a blueprint.

The scope of this step is a week consisting of three ORT sessions.
Because all three sessions serve the same specialty, the specialty assignment
variables $V_{hq}$ and mismatch penalty $D_{hq}$ from the Step~3 model are
no longer needed and are dropped.
Furthermore, since the problem scope is defined by exactly three ORT sessions,
all sessions are always open:
\begin{equation}
  Z_h = 1 \qquad \forall\, h \in H.
  \label{eq:allopen}
\end{equation}

\subsection{Patient Subgroup Definition}
\label{sec:clustering}

Because the blueprint must be constructed before the waiting list is revealed,
we cannot decide on individual patient level.
In addition, the patient mix consists of a large number of different surgery
types that only sporadically occur in the waiting list.
Therefore, patients are aggregated into higher-level subgroups $g \in G$ by
clustering procedure types based on the empirical mean and standard deviation
of surgery duration using $k$-means clustering in the $(\mu, \sigma)$ space.
This data-driven approach extends the label-based three-group classification
of Sheweita et al.~\cite{sheweita2021} by capturing scheduling risk ($\sigma$)
alongside expected OR consumption ($\mu$).
Three clusters are identified:
\textbf{S~(Short)}: low mean, low variance;
\textbf{M~(Medium)}: moderate mean and variance; and
\textbf{L~(Long)}: high mean and/or variance,
giving $G = \{\mathrm{S}, \mathrm{M}, \mathrm{L}\}$.
Within each subgroup, durations are approximately exchangeable, justifying the
subgroup count $N_{gh}$ as the relevant decision variable at Stage~1.

\subsection{Three-Stage Structure and Scenario Tree}
\label{sec:three_stage_structure}

Uncertainty resolves in two phases.
\textbf{Stage~1}~(MSSP): the blueprint $N_{gh}$ is fixed before any waiting
list information is known.
\textbf{Stage~2}~(SSP): the waiting list realisation $r \in R$ is revealed;
patients are assigned, sequenced, and given appointment times using variables
$X_{ijh}$, $Y_{ph}$, $A_p$, $Z_h$, $U_{ph}$, now indexed by~$r$.
\textbf{Stage~3}~(realisation): surgery durations are observed and the
schedule is evaluated using variables $S_{ps}$, $W_{ps}$, $I_{pp's}$,
$O_{hs}$, now indexed by~$rs$.

A node-oriented formulation is adopted, indexing variables by nodes in a
scenario tree $\mathcal{T}$ with root $n_0$, waiting-list nodes $r \in R$
with realised sets of patients, and leaf nodes $rs$ with patients and
treatment durations.
Non-anticipativity is enforced structurally: $N_{gh}$ is defined at $n_0$
and is shared across all branches; Stage~2 variables are defined at node~$r$
and are shared across all duration sub-branches $s \in S$; Stage~3 variables
are $rs$-specific.
No explicit non-anticipativity constraints are therefore needed.
In a time-oriented formulation the same requirement would demand explicit
equality constraints (i.e.\ restricting the first-stage blueprint to be
consistent: $N_{ghr} = N_{ghr'}$ for all $r, r' \in R$), which the
node-oriented formulation does not need.

\subsection{Extended Model}
\label{sec:extended_model}

The Step~3 model is extended as follows.
The set $r \in R$ indexes waiting-list realisations, analogously to how
$s \in S$ indexes duration scenarios; in total the scenario tree contains
$|R| \times |S|$ leaves.
Each realisation $r$ defines a patient set $P_r$ with subgroup partition
$P_{gr}$ for all $g \in G$.
The probabilities $\pi_r$ and $\pi_{s|r}$ satisfy $\sum_{r} \pi_r = 1$ and
$\sum_{s} \pi_{s|r} = 1$.
The new Stage~1 variable is $N_{gh} \in \mathbb{Z}_{\geq 0}$, the number of
patients of subgroup~$g$ that may be allocated to session~$h$ under the
blueprint.

\paragraph{Objective function.}
Since all sessions are ORT, the specialty mismatch penalty is dropped.
To prevent the blueprint from becoming empty, without a cost on $N_{gh}$
the optimiser would set every quota to $\max_r |P_{gr}|$, imposing no
planning structure, a penalty $\beta_5$ is added for each blueprint slot.
This creates tension between blueprint specificity and operational flexibility:
a tighter blueprint (small $N_{gh}$) reduces planning ambiguity but may
constrain Stage~2 unnecessarily, while a looser blueprint provides more
flexibility at a higher planning cost.
The parameter $\beta_5$ was set to~$1$, ensuring the blueprint remains
structured without trivially allocating all capacity; a sensitivity analysis
on $\beta_5$ could further improve solution quality but is left for future work.

\begin{align}
  \min \; & \beta_5 \sum_{h \in H} \sum_{g \in G} N_{gh}
    + \sum_{r \in R} \pi_r \sum_{s \in S} \pi_{s|r} \left(
        \beta_1 \sum_{p} W_{prs} \right. \nonumber \\
  & \left. \quad
      + \; \beta_2 \sum_{\substack{p,p' \in P_r\\p \neq p'}} I_{pp'rs}
      + \beta_3 \sum_{h} O_{hrs}
    \right) \label{eq:three_stage_obj}
\end{align}

\paragraph{Blueprint linking constraint.}
The blueprint links Stage~1 to Stage~2 by capping the number of patients of
each subgroup per session:

\begin{equation}
  \sum_{p \in P_{gr}} Y_{phr} \leq N_{gh}
  \qquad \forall\, g \in G,\; h \in H,\; r \in R
  \label{eq:blueprint_link}
\end{equation}

\noindent The inequality (rather than equality) allows under-utilisation of
blueprint capacity: when realization~$r$ contains fewer patients of subgroup~$g$
than the quota, the constraint is non-binding, preserving feasibility.
Conversely, when patients are abundant the constraint is active and restricts
Stage~2 assignments to at most $N_{gh}$ patients per group per session.

All remaining Step~3 constraints (excluding the specialty-related
constraints~\eqref{eq:spec} and~\eqref{eq:viol}) carry forward as Stage~2
constraints (indexed by~$r$) and Stage~3 constraints (indexed by~$rs$).
The deterministic equivalent scales as
$\mathcal{O}(|R| \cdot |P|^2 \cdot |H|)$ in binary variables and
$\mathcal{O}(|R| \cdot |S| \cdot (|P|^2 + |H|))$ in continuous variables;
tractability is maintained by applying SAA with moderate $|R|$ and~$|S|$.

\paragraph{Session symmetry-breaking.}
The three sessions are identical (equal opening and closing times, no
session-specific data), so any permutation of the session indices yields an
equivalent solution. To remove this $|H|!$ symmetry, which would otherwise
inflate the branch-and-bound tree, the blueprint columns are ordered by a
lexicographic capacity key that weights L above M above S:
$100\,N_{\mathrm{L}h} + 10\,N_{\mathrm{M}h} + N_{\mathrm{S}h}$ is required to be
non-increasing in~$h$. Because the sessions are genuinely exchangeable this cut
preserves optimality while fixing a canonical labelling (session~$0$ carries the
heaviest case mix).

\subsection{Instance Design}
\label{sec:bp_instance}

Surgery types are clustered on the full ORT history
(\texttt{ort\_patient\_data.csv}); the resulting three subgroups are summarised
in Table~\ref{tab:bp_clusters}. The clustering is dominated by short procedures
(50 of 68 types), with a small number of long, high-variance types forming the
L~group. Each waiting-list realisation is generated by drawing three documented
ORT sessions at random from \texttt{step\_5\_sample\_sessions.csv} and pooling
their patients; we use $|R| = 3$ realisations with equal probability
$\pi_r = 1/3$, and $|S| = 5$ duration scenarios per realisation drawn by
importance sampling ($k = 1$). The blueprint penalty is set to $\beta_5 = 1$.
The three sampled realisations contain $16$, $14$ and $14$ patients with
subgroup compositions $(15\mathrm{S},1\mathrm{M},0\mathrm{L})$,
$(12\mathrm{S},1\mathrm{M},1\mathrm{L})$ and
$(12\mathrm{S},2\mathrm{M},0\mathrm{L})$, respectively.

\begin{table}[h]
  \centering
  \caption{Surgery-type subgroups from $k$-means clustering on the ORT history.}
  \label{tab:bp_clusters}
  \begin{tabular}{lrrr}
    \toprule
    Subgroup & \#\,types & mean duration (min) & mean $\sigma$ (min) \\
    \midrule
    S (Short)  & 50 &  63.5 & 17.1 \\
    M (Medium) & 15 & 126.3 & 29.6 \\
    L (Long)   &  3 & 247.5 & 77.7 \\
    \bottomrule
  \end{tabular}
\end{table}

\subsection{Blueprint Results and Managerial Insights}
\label{sec:bp_results}

The resulting blueprint is reported in Table~\ref{tab:bp_quota}. Two structural
features stand out. First, the \emph{per-subgroup totals} (15~S, 2~M, 1~L) equal
the largest subgroup count observed across the three realisations; with
$\beta_5 = 1$ the model provisions exactly enough capacity to absorb the
worst-case realisation of each subgroup and no more. This part of the solution
is provably tight: the lower bound returned by the solver equals the slot cost,
confirming that no smaller blueprint can remain feasible. Second, the single
long case is isolated in one session while the medium cases are concentrated in
another, leaving the remaining capacity for short, high-throughput procedures.

\begin{table}[h]
  \centering
  \caption{Blueprint quota $N_{gh}$ (patients of subgroup $g$ allowed in session $h$).}
  \label{tab:bp_quota}
  \begin{tabular}{lrrrr}
    \toprule
    Subgroup & Session 0 & Session 1 & Session 2 & Total \\
    \midrule
    S (Short)  & 6 & 3 & 6 & 15 \\
    M (Medium) & 0 & 2 & 0 & 2  \\
    L (Long)   & 1 & 0 & 0 & 1  \\
    \midrule
    Total      & 7 & 5 & 6 & 18 \\
    \bottomrule
  \end{tabular}
\end{table}

The managerial reading is that a robust weekly blueprint should reserve one
session for the long, high-variance procedures and dedicate the others to a
high volume of short cases, rather than spreading long cases evenly. Isolating
the variance-driving cases limits how far a single overrun can cascade into
downstream waiting and overtime, while the short-case sessions remain dense and
predictable.

\paragraph{Tractability.}
As in Step~3, the embedded sequencing constraints carry large Big-$M$
coefficients, and here this weakness is compounded by the three nested Stage-2
scheduling problems. Within a $600$-second time limit the solver therefore
returns a high-quality feasible blueprint with a sizeable residual optimality
gap on the \emph{operational} cost component, while the blueprint capacities
themselves are at the proven lower bound. Reducing $|S|$, using random rather
than importance sampling for a more stable cost estimate, or strengthening the
second-stage relaxation would tighten this gap; a full $\beta_5$ sensitivity
analysis is left for future work.

\subsection{Blueprint Stability Across Sampling Seeds}
\label{sec:bp_consistency}

Because the blueprint is constructed from a finite Sample Average Approximation
(SAA) of the underlying uncertainty---$|R|$ waiting-list realisations and
$|S|$ duration scenarios per realisation---a single solve reflects one
particular random draw. To establish that the recommended quota table
$N_{gh}$ is a property of the data rather than an artefact of one sample, we
re-solve the three-stage program over a sequence of independent seeds
$\theta \in \{42, 43, 44, \dots\}$ and compare the resulting blueprints.

For the comparison to be meaningful the subgroup definitions must be held
invariant across seeds. The $k$-means clustering of
Section~\ref{sec:clustering} depends on a (random) initialisation, so re-seeding
it together with the sampler would permute the S/M/L labels and render the
quota tables incommensurable. We therefore decouple the two sources of
randomness: the clustering is computed \emph{once} on the full ORT history with
a fixed clustering seed, while the seed $\theta$ governs only the operational
draws---which three documented sessions form each realisation
(Section~\ref{sec:bp_instance}) and the duration scenarios sampled within each
realisation. All structural parameters ($|R|$, $|S|$, $\beta_5$, and the
sampling method) are likewise fixed.

Each seed therefore produces a blueprint over the \emph{same} subgroup
definitions, and we record for every run its settings, its objective
decomposed into the slot cost
$\beta_5\sum_{g,h} N_{gh}$ and the residual operational cost, the attained
optimality gap, and the complete quota table $N_{gh}$, one row per seed in a
cumulative results file. Stability is then read directly from this table:
agreement of the per-subgroup totals $\sum_h N_{gh}$ across seeds indicates that
the provisioned capacity is robust to the particular waiting lists sampled,
while agreement of the per-session allocation indicates that the qualitative
managerial recommendation---isolating the long, high-variance cases and
densifying the short-case sessions (Section~\ref{sec:bp_results})---is itself a
stable feature rather than sampling noise.

\subsection{Operational Cost of Imposing the Blueprint}
\label{sec:bp_apply}

The blueprint is a tactical template fixed before the waiting list is known;
its purpose is to constrain the downstream operational schedule. To quantify the
operational price of obeying it, we evaluate a \emph{given} blueprint as an
a~priori restriction on an independent scheduling instance and compare the
resulting schedule against the unrestricted optimum.

\paragraph{Comparison instance.}
A dedicated instance is built that is structurally identical to the blueprint's
Stage~2: three documented ORT sessions are drawn at random from the ORT history
and their patients pooled, mirroring the realisation construction of
Section~\ref{sec:bp_instance} but retaining the recorded session assignment and
intra-session sequence so that the operational Step~1--3 chain applies
unchanged. The instance is fixed once and reused across both solves so that the
two differ only in the presence of the blueprint. Its surgery types are mapped
to the subgroups $g \in G$ using the same fixed clustering, so the quotas
$N_{gh}$ are directly applicable.

\paragraph{Restricted versus unrestricted schedules.}
The operational Step~1--3 warm-start chain (timing only; sequencing and timing;
full assignment) is solved twice on this instance.
In the \emph{baseline} run the chain is solved as usual, leaving the
patient-to-session assignment free in the full step.
In the \emph{restricted} run the blueprint linking
constraint~\eqref{eq:blueprint_link} is imposed on the (now fixed) capacities
$N_{gh}$,
\begin{equation}
  \sum_{p \in P_g} Y_{ph} \leq N_{gh}
  \qquad \forall\, g \in G,\; h \in H,
  \label{eq:bp_apply_link}
\end{equation}
where $P_g$ is the subgroup partition of the instance.
The restriction is binding only in the full step: Steps~1 and~2 fix the
assignment $Y_{ph}$ from the documented schedule, so they are identical across
the two runs, whereas the full step optimises $Y_{ph}$ freely and is therefore
the locus at which the blueprint reshapes the solution.

\paragraph{Evaluation.}
Both full-step policies are evaluated out of sample on a single, independently
sampled scenario set, following the same shared-evaluation protocol used for the
SAA convergence and VSS analyses, which removes evaluation-sample noise from the
comparison. The quantity of interest is the difference in out-of-sample expected
cost between the restricted and baseline policies: the \emph{operational cost of
the blueprint}, i.e.\ the flexibility forgone by committing to the tactical
template. Since the restricted feasible region is a subset of the baseline one,
the restricted in-sample objective is an upper bound on the baseline objective,
so this difference is non-negative up to the residual optimality gaps of the two
solves. Reporting it alongside the realised per-session subgroup composition of
each schedule shows both \emph{how much} structure costs and \emph{where} the
blueprint forces the assignment to differ from the unconstrained optimum.